\documentclass[12pt,letterpaper]{article}

\usepackage[margin=0.75in]{geometry}
\usepackage{amsmath,amssymb}
\usepackage{enumitem}

\pagestyle{empty}
\setlength{\parindent}{0pt}
\setlength{\parskip}{0.5em}
\setlength{\emergencystretch}{2em}

\begin{document}

\begin{center}
  {\large\bfseries MATH 15200: Calculus II}\\
  {\large\bfseries Homework 1: Limits and Continuity}
\end{center}

\vspace{0.25em}
\hrule
\vspace{0.8em}

\textbf{Due:} Friday, October 9, 2026, at 11:59 PM.

This is the written portion of Homework 1; complete the routine online portion
separately in Edfinity. Upload your written solutions to Gradescope, via Canvas. You may
collaborate, but you must write and submit your own solutions independently.

\vspace{0.3em}

\noindent\fbox{%
  \parbox{0.955\textwidth}{%
    \textbf{Definition.}
    Let $f$ be a real-valued function defined near $a$, possibly not at $a$.
    We say
    \[
      \lim_{x\to a} f(x)=L
    \]
    if, for every $\varepsilon>0$, there exists $\delta>0$ such that, whenever
    $x$ is in the domain of $f$,
    \[
      0<|x-a|<\delta
      \quad\Longrightarrow\quad
      |f(x)-L|<\varepsilon.
    \]
    If $f(a)$ is defined, then $f$ is \textbf{continuous at $a$} when
    \[
      \lim_{x\to a} f(x)=f(a).
    \]
  }%
}

\vspace{0.8em}

Rewrite each of the problems below in terms of the epsilon-delta definiton above. 
Explicitly state your choice of \(\delta\) when necessary.

\begin{enumerate}[label=\textbf{\arabic*.},leftmargin=2em,itemsep=1.25em]
  \item Prove directly from the epsilon--delta definition that
    \[
      \lim_{x\to 2}\bigl(x^2-1\bigr)=3.
    \]

    \emph{Hint:} Impose a simple bound such as $|x-2|<1$.

  \item Define a function \(f\) by the piecewise expression
    \[
      f(x)=
      \begin{cases}
        0, & x<0,\\
        1, & x\geq 0.
      \end{cases}
    \]
    Prove directly from the definition that $f$ is not continuous at $0$.

    \emph{Hint:} Choose a fixed positive $\varepsilon$ and consider points
    immediately to the left of $0$.

  \item Suppose that $f$ and $g$ are real-valued functions defined near $a$ and
    that
    \[
      \lim_{x\to a}f(x)=L
      \qquad\text{and}\qquad
      \lim_{x\to a}g(x)=M.
    \]
    Prove from the epsilon--delta definition that
    \[
      \lim_{x\to a}\bigl(f(x)+g(x)\bigr)=L+M.
    \]

    \emph{Hint:} Apply the two hypotheses using $\varepsilon/2$ and combine
    the resulting choices of $\delta$.
\end{enumerate}

\end{document}
